\documentclass{practice}

\title{11}
\date{15.04.26}

\usepackage{crysymb}

\begin{document}
\maketitle

\begin{task}{PKCS\#11}
  PKCS\#11 (or \emph{Cryptoki}) defines an interface for manipulating cryptographic tokens and is often used by HSMs and smart cards.
  The tokens are typically X.509 certificates or public/private keys, but can also be symmetric keys.
  Depending on the device's capabilities, the Cryptoki interface can be used to use, extract, create, delete, and modify those objects.

  Likely the most common PKCS\#11 module is the one provided by \href{https://github.com/opensc/opensc}{\emph{OpenSC}}.
  Many vendors will have their own (e.g. Yubico's YKCS11) which allows for easier manipulation/discovery of objects on the vendor's device.

  The OpenSC project also provides the \texttt{pkcs11-tool} utility, which can be used to interact with cryptographic tokens.
  For this task, install the utility (or the full OpenSC toolset, which also contains it).
\end{task}

\begin{task}{ID card's PKCS\#11}
  Most of you have (or should have) and Estonian ID card or residence permit card.
  If you have neither, but the ID card from your country has a chip, it will likely also have a PKCS\#11 interface.

  For this task, we will use \texttt{pkcs11-tool} to play around with the ID card.

  To list token slots identified by the tool:
  \begin{cmdblock}
pkcs11-tool -L
  \end{cmdblock}
  \begin{figure}[h]
    \includegraphics[width=\textwidth]{../../images/pkcs11-id-list}
  \end{figure}
  Notice that the tokens in the slots have \enquote{readonly} set, meaning that we cannot edit those tokens.
  This is good, since otherwise people could accidentally overwrite their eID keypair.
  \enquote{login required} means that the PIN must be provided in order to interact with private objects.

  \newpage

  To list cryptographic mechanisms supported by the card:
  \begin{cmdblock}
pkcs11-tool -M
  \end{cmdblock}
  \begin{figure}[h]
    \includegraphics[width=\textwidth]{../../images/pkcs11-id-mechanisms}
  \end{figure}
  We can see that my ID card does not support RSA.

  Let us try to see what tokens are actually on the card:
  \begin{cmdblock}
pkcs11-tool -O
  \end{cmdblock}
  \begin{figure}[h]
    \includegraphics[width=\textwidth]{../../images/pkcs11-id-objects}
  \end{figure}
  Some objects seem to be missing: we know that there should be a private key somewhere, and right now only the authentication objects are displayed.

  To display also the private key object, we must \enquote{log in} to the card.
  \begin{cmdblock}
pkcs11-tool -O --login
  \end{cmdblock}
  \begin{figure}[h]
    \includegraphics[width=\textwidth]{../../images/pkcs11-id-objects-login}
  \end{figure}
  We can now see the private key object, but still not the objects tied to the signing certificate.

  Notice that \enquote{sensitive} and \enquote{always sensitive} have been set, which means that the key has always been protected, i.e. it cannot be viewed or exported as plaintext.
  This typically also means that it was generated on the device.

  A sensitive key might still be exportable if it is wrapped; i.e. the encrypted private key can be exported, for example, to re-import it into a different HSM or smart card.
  Since \enquote{never extractable} is set, the token cannot, under any circumstance, be exported from the device.

  To see the signing certificate objects, we must specify the slot since by default the tool selects the first slot (in this case slot 0).
  \begin{cmdblock}
pkcs11-tool -O --login --slot 1
  \end{cmdblock}

  Let us fetch the signing certificate.
  \begin{cmdblock}
pkcs11-tool --slot 1 -r --type cert \
  --label Allkirjastamine > cert.der
  \end{cmdblock}
  You can verify that it worked by using
  \begin{cmdblock}
openssl x509 -in cert.der -noout -subject
  \end{cmdblock}
  \begin{figure}[h]
    \includegraphics[width=\textwidth]{../../images/pkcs11-id-sigcert}
  \end{figure}

  Finally, let us sign some data with the ID card.
  For fun, let us use the Authentication key, since DigiDoc only allows you to sign data using the Signature key.
  \begin{cmdblock}
echo "Message" > message.txt

pkcs11-tool --sign -m ECDSA-SHA384 --slot 0 \
  --input-file message.txt --output-file message.sig
  \end{cmdblock}
  ECDSA-SHA384 means that the device's own functionality will be used to hash the data with SHA-384, and that hash is then signed with ECDSA.

  We can also use the device to verify the signature:
  \begin{cmdblock}
pkcs11-tool --verify -m ECDSA-SHA384 --slot 0 \
  --input-file message.txt --signature-file message.sig
  \end{cmdblock}
  \begin{figure}[h]
    \includegraphics[width=\textwidth]{../../images/pkcs11-id-signature}
  \end{figure}

  Let us try to verify the signature with OpenSSL.
  \begin{cmdblock}
pkcs11-tool --slot 0 -r --type cert \
  --label Isikutuvastus > cert.der

openssl x509 -in cert.der -pubkey -noout > pub.pem

openssl pkeyutl -verify -digest sha384 \
  -in message.txt -pubin -inkey pub.pem -sigfile message.sig
  \end{cmdblock}
  which results in \texttt{Signature Verification Failure}.

  The issue is that the device exports the ECDSA signature as a raw concatenation of the r and s values, whereas OpenSSL expects the signature to be ASN.1-encoded according to \href{https://datatracker.ietf.org/doc/html/rfc3279#section-2.2.3}{RFC 3279}.
  \begin{Verbatim}
Ecdsa-Sig-Value  ::=  SEQUENCE  {
           r     INTEGER,
           s     INTEGER  }
  \end{Verbatim}

  You can do the conversion using the helper script \texttt{wrap\_ecdsa.py} (which requires the \texttt{pyasn1} Python package) on the course GitHub.
  \begin{cmdblock}
python3 wrap_ecdsa.py message.sig wrapped.sig

openssl pkeyutl -verify -digest sha384 \
  -in message.txt -pubin -inkey pub.pem -sigfile wrapped.sig
  \end{cmdblock}
  which should result in \texttt{Signature Verified Successfully}.
\end{task}

\begin{task}{EstEID}
  The Estonian eID application on your card is documented in one of the following:
  \begin{itemize}
    \item \href{https://www.id.ee/wp-content/uploads/2026/03/id1developerguide2025.pdf}{\textit{IDEMIA EstEID developer guide}} for cards issued until 14.11.2025
    \item \href{https://www.id.ee/wp-content/uploads/2025/11/ged24035_015_tdc_est_eid_developer_guide.pdf}{\textit{Thales EstEID developer guide}} for cards issued starting from 17.11.2025
  \end{itemize}

  See \S 9.2.1 (IDEMIA) or \S 5.3.11.1 (Thales) of the documentation to see what personal information is stored on your card's eID application.
  This information is e.g. what is (partly) read by terminals when you use your ID card as a loyalty card.

  \begin{tcolorbox}[title=Note]
    You can only do the practical part of this task if your ID card was issued before 17.11.2025.
    Since 17.11.2025, cards have been manufactured by Thales rather than IDEMIA, and I have not been able to test my script with an actual ID card, as I do not trust the test card sufficiently.
  \end{tcolorbox}

  Use the script \texttt{esteid\_info.py} (which requires the \texttt{pyscard} Python package) to view the personal information stored on your ID card.
\end{task}
\end{document}
